\documentclass{article}
\usepackage{amsmath}
\usepackage{amssymb}
\usepackage{geometry}
\geometry{margin=1in}

\title{Modelo de Regresion Econometrica: Adopcion vs. Fallos en Produccion}
\author{Sistema Multi-Agente de Investigacion}
\date{Abril 2026}

\begin{document}

\maketitle

\section{Especificacion del Modelo}

Proponemos una regresion econometrica para medir la relacion entre tasa de adopcion de IA y tasa de fallos en produccion:

\begin{equation}
\ln(F_t) = \beta_0 + \beta_1 \ln(A_{t-1}) + \beta_2 \ln(D_{t-1}) + \beta_3 T_t + \beta_4 Q_t + \epsilon_t
\end{equation}

Donde:
\begin{itemize}
    \item $F_t$ = Tasa de fallos en produccion en tiempo $t$ (bugs criticos por 1000 LOC)
    \item $A_{t-1}$ = Tasa de adopcion de IA en $t-1$ (porcentaje de codigo generado por IA)
    \item $D_{t-1}$ = Deuda tecnica acumulada en $t-1$ (escala logaritmica)
    \item $T_t$ = Inversion en training/quality assurance en $t$ (horas por developer)
    \item $Q_t$ = Calidad de la arquitectura en $t$ (indice 0-1)
    \item $\epsilon_t$ = Termino de error
\end{itemize}

\section{Hipotesis sobre los Coeficientes}

\begin{table}[h]
\centering
\begin{tabular}{|c|c|c|l|}
\hline
\textbf{Coef} & \textbf{Signo Esperado} & \textbf{Magnitud} & \textbf{Interpretacion} \\
\hline
$\beta_1$ & $+$ & 0.3 - 0.7 & Mayor adopcion sin control $\rightarrow$ mas fallos \\
$\beta_2$ & $+$ & 0.4 - 0.8 & Mayor deuda tecnica $\rightarrow$ mas fallos \\
$\beta_3$ & $-$ & -0.2 - -0.5 & Mas training $\rightarrow$ menos fallos \\
$\beta_4$ & $-$ & -0.5 - -1.0 & Mejor arquitectura $\rightarrow$ menos fallos \\
\hline
\end{tabular}
\end{table}

\section{Modelo Dinamico de Transicion}

Sistema de ecuaciones diferenciales para la evolucion del mercado:

\begin{align}
\frac{dS}{dt} &= r_S \cdot S \cdot \left(1 - \frac{S + P}{K}\right) - \delta_{SP} \cdot S \cdot P \\
\frac{dP}{dt} &= r_P \cdot P \cdot \left(1 - \frac{S + P}{K}\right) + \delta_{SP} \cdot S \cdot P - \mu \cdot D \\
\frac{dD}{dt} &= \alpha \cdot P - \beta \cdot D
\end{align}

Donde:
\begin{itemize}
    \item $S(t)$ = Empresas de software tradicional (millones)
    \item $P(t)$ = Empresas de parches/remediacion (millones)
    \item $D(t)$ = Deuda tecnica total del mercado (billones USD)
    \item $K$ = Capacidad de mercado (constante de carrying capacity)
    \item $r_S$ = Tasa intrinseca de crecimiento de empresas tradicionales
    \item $r_P$ = Tasa intrinseca de crecimiento de empresas de remediacion
    \item $\delta_{SP}$ = Tasa de transicion de tradicional a remediacion
    \item $\mu$ = Tasa de mortalidad de empresas por deuda tecnica
    \item $\alpha$ = Tasa de generacion de deuda por vibe coding
    \item $\beta$ = Tasa de remediacion de deuda
\end{itemize}

\section{Parametros Estimados}

Basado en datos de mercado 2024-2026:

\begin{table}[h]
\centering
\begin{tabular}{|c|c|c|}
\hline
\textbf{Parametro} & \textbf{Valor} & \textbf{Fuente} \\
\hline
$r_S$ & 0.15 & CAGR empresas software tradicional \\
$r_P$ & 0.35 & CAGR empresas ciberseguridad/remediacion \\
$\delta_{SP}$ & 0.05 & Tasa de conversion observada \\
$\mu$ & 0.08 & Tasa de fallo empresas con alta deuda \\
$\alpha$ & 2.5 & Generacion de deuda por vibe coding \\
$\beta$ & 1.2 & Capacidad de remediacion del mercado \\
$K$ & 100 & Capacidad de mercado (normalizado) \\
\hline
\end{tabular}
\end{table}

\section{Puntos de Equilibrio}

Los puntos de equilibrio del sistema se encuentran resolviendo:

\begin{align}
0 &= r_S \cdot S^* \cdot \left(1 - \frac{S^* + P^*}{K}\right) - \delta_{SP} \cdot S^* \cdot P^* \\
0 &= r_P \cdot P^* \cdot \left(1 - \frac{S^* + P^*}{K}\right) + \delta_{SP} \cdot S^* \cdot P^* - \mu \cdot D^* \\
0 &= \alpha \cdot P^* - \beta \cdot D^*
\end{align}

\subsection{Equilibrio 1: Dominancia Tradicional}

$(S^*, P^*, D^*) = (K, 0, 0)$

Este equilibrio es inestable si $\delta_{SP} > 0$.

\subsection{Equilibrio 2: Estado Estacionario Mixto}

Resolviendo el sistema no lineal:

\begin{align}
S^* &= \frac{\beta K (r_P - \mu \alpha / \beta)}{\beta r_P + \alpha \delta_{SP} K} \\
P^* &= \frac{\alpha K (r_S - \delta_{SP} P^*)}{\beta r_S + \alpha \delta_{SP} K} \\
D^* &= \frac{\alpha}{\beta} P^*
\end{align}

\section{Simulacion Numerica}

Condiciones iniciales (2024):
\begin{itemize}
    \item $S_0 = 85$ (85\% empresas tradicionales)
    \item $P_0 = 15$ (15\% empresas de remediacion)
    \item $D_0 = 2.1$ (deuda tecnica acumulada)
\end{itemize}

Proyeccion 2030:
\begin{itemize}
    \item $S_{2030} \approx 45$ (45\%)
    \item $P_{2030} \approx 40$ (40\%)
    \item $D_{2030} \approx 5.2$ (deuda tecnica)
\end{itemize}

\section{Implicaciones de Politica}

\subsection{Condicion para Estabilidad}

Para evitar la explosion de deuda tecnica:

\begin{equation}
\beta > \alpha \cdot \frac{P}{D_{critica}}
\end{equation}

Esto implica que la capacidad de remediacion debe crecer mas rapido que la generacion de deuda.

\subsection{Inversion Optima en Remediacion}

La inversion optima $I^*$ en capacidad de remediacion satisface:

\begin{equation}
\frac{\partial \beta(I)}{\partial I} \cdot D = \frac{\partial C(I)}{\partial I}
\end{equation}

Donde $C(I)$ es el costo de la inversion.

\section{Extension: Modelo Estocastico}

Para capturar incertidumbre, agregamos terminos estocasticos:

\begin{equation}
d\mathbf{X} = \mathbf{f}(\mathbf{X}) dt + \mathbf{g}(\mathbf{X}) d\mathbf{W}
\end{equation}

Donde:
\begin{itemize}
    \item $\mathbf{X} = [S, P, D]^T$
    \item $\mathbf{f}(\mathbf{X})$ = vector de dinamicas deterministicas
    \item $\mathbf{g}(\mathbf{X})$ = matriz de difusion
    \item $d\mathbf{W}$ = proceso de Wiener multidimensional
\end{itemize}

\section{Conclusion}

El modelo predice una transicion inevitable hacia un mercado donde las empresas de remediacion representan un porcentaje significativo (40\%+) del ecosistema de software para 2030, impulsada por la acumulacion de deuda tecnica generada por el vibe coding.

La condicion critica para la sostenibilidad del ecosistema es que $\beta > \alpha$, es decir, que la capacidad de remediacion crezca mas rapido que la generacion de deuda tecnica.

\end{document}
